\documentclass[../../../uem_mei_q.tex]{subfiles}

\begin{document}
    次の空欄中\textsf{ア}，\textsf{ウ}～\textsf{ク}に当てはまるものを選択肢の中から選び%
    その記号をマークせよ．空欄\textsf{イ}には最も適切なものを選択肢の中から選び%
    その記号をマークせよ．それ以外の空欄に当てはまる0から9までの数字を解答用紙の所定の欄に%
    マークせよ．

    \vspace{.5\zw}
    $n$を2以上の自然数とする．原点がOである座標平面を考え，$C$を原点中心で半径が1の円とする．%
    \vspace{1pt}%
    この平面上の点AとPの座標をそれぞれ$(0,\;1)$と$\left(\myfracb{1}{n},\;0\right)$とおく．%
    点Qは次の条件をすべて満たすような点であるとする．
    
    \begin{myenua}
        \item 点Qは円$C$上にある．
            \label{uem_mei_2025_uni1_03_q:1}
        \item $\overrightarrow{\mathrm{AQ}}\cdot\overrightarrow{\mathrm{PQ}}=0$
            \label{uem_mei_2025_uni1_03_q:2}
        \item 点QとAの座標は異なる．
            \label{uem_mei_2025_uni1_03_q:3}
    \end{myenua}
    以下の問いに答えよ．
    \\    
    \begin{minipage}{\linewidth}
        \vspace{.5\baselineskip}
        \centering
            \includegraphics%
            {\subfix{fig_uem_mei_2025_uni1_03/fig_uem_mei_2025_uni1_03_01_q.pdf}}
    \end{minipage}
    \vspace{.1\baselineskip}

    \begin{enumerate}
        \item 線分APの長さは $\myfraca{\sqrt{\rule[0pt]{0pt}{13pt}%
            \text{\myboxd{　ア　}}}}{n}$ である．
        \item 点Qの座標を$(x,\;y)$とおいて，条件\ref{uem_mei_2025_uni1_03_q:1}，%
            \ref{uem_mei_2025_uni1_03_q:2}を$n$，$x$，$y$の式で表すことにより，%
            $\text{\myboxd{　イ　}}=0$を得る．さらに条件\ref{uem_mei_2025_uni1_03_q:3}%
            を考慮することにより，点Qの座標は%
            $\left(\myfraca{\text{\myboxd{　ウ　}}}{\text{\myboxa{　ア　}}},%
            \;\myfraca{\text{\myboxd{　エ　}}}{\text{\myboxa{　ア　}}}\right)$%
            であることがわかる．
        \item 3つの点O，P，Aを通る円を$C_1$とすると，$C_1$は点Qも通る．%
            このことに注意すると，$\sin\angle\mathrm{GAP}=%
            \myfraca{\text{\myboxd{　オ　}}}{\text{\myboxc{　ア　}}}$ である．
        \item 三角形PAQについて\\
            $\mathrm{AP}:\mathrm{AQ}:\mathrm{PQ}=%
            \,\text{\myboxc{　ア　}}\,:\,\text{\myboxd{　カ　}}\,:\,\text{\myboxd{　キ　}}$ \:%
            である．
        \item 三角形PAQの内接円の半径を$r$とすると，\\
            $\mathrm{AP}:r=\,\text{\myboxc{　ア　}}\,:\,\text{\myboxd{　ク　}}$ \:である．
        \newpage
        \item $n=2$ とする．点Iを三角形PAQの内接円の中心であるとしたとき
            \begingroup
                \setlength{\abovedisplayskip}{7pt}
                \setlength{\belowdisplayskip}{5pt}
                \setlength{\jot}{0pt}
                \begin{align*}
                    \overrightarrow{\mathrm{QI}}=%
                    \myfraca{1}{\text{\myboxb{　ケ　}}}\overrightarrow{\mathrm{QA}}+%
                    \myfraca{1}{\text{\myboxb{　コ　}}}\overrightarrow{\mathrm{QP}}
                \end{align*}
            \endgroup
            であり，$\overrightarrow{\mathrm{OI}}=%
            \overrightarrow{\mathrm{OQ}}+\overrightarrow{\mathrm{QI}}$ %
            であることから，点Iの座標は%
            $\biggl(\myfraca{1}{\text{\myboxb{　サ　}}},\;\myfraca{1}{\text{\myboxb{　シ　}}}\biggr)$%
            である．
    
        \vspace{.5\baselineskip}
        \textsf{ア，ウ～クの解答群}
        
        　\begin{tabularx}{\dimexpr\linewidth-2\zw}{@{}LLLL@{}}
            ⓪\; $0$ & ①\; $1$ & ②\; $n$ & ③\; $2n$ \\ [2pt]
            ④\; $n^2$ & ⑤\; $2n^2$ & ⑥\; $n+1$ & ⑦\; $n-1$ \\ [2pt]
            ⑧\; $n^2+1$ & ⑨\; $n^2-1$
        \end{tabularx}

        \vspace{.5\baselineskip}
        \textsf{イの解答群}
        
        \vspace{3pt}
        　\begin{tabularx}{\dimexpr\linewidth-2\zw}{@{}LLLL@{}}
            ⓪\; $x+\myfracc{y}{n}+1$ & ①\; $x-\myfracc{y}{n}+1$ & ②\; $x+\myfracc{y}{n}-1$ \\[9pt]
            ③\; $x-\myfracc{y}{n}-1$ & ④\; $\myfracb{x}{n}+y+1$ & ⑤\; $\myfracb{x}{n}-y+1$ \\[9pt]
            ⑥\; $\myfracb{x}{n}+y-1$ & ⑦\; $\myfracb{x}{n}-y-1$ & ⑧\; $x-ny+1$ \\[7pt]
            ⑨\; $x+ny+1$
        \end{tabularx}
    \end{enumerate}
\end{document}
